\ifajar
  \section{Latihan}
\else
  \section{Bab 10: Identity, Authentication, dan Authorization}
\fi

Seluruh file yang dipakai pada latihan ini tersedia di
\berkas{sources/chapter10/}, dihitung relatif terhadap akar repositori.
Latihan ini hanya membutuhkan Python 3, tanpa Docker dan tanpa basis data.

Seluruh permintaan diarahkan ke aplikasi yang berjalan di komputer sendiri.
Kata sandi pada \berkas{sources/chapter10/pengguna.json} hanya untuk latihan,
dan kunci tandanya dibaca dari \textit{environment} \texttt{KUNCI\_TOKEN}.
Hasil tiap mahasiswa akan sama pada bagian status, dan berbeda pada bagian
penanda yang memang acak.

Penyiapannya satu perintah, seperti pada Listing \ref{lst:c10-siapkan}.

\begin{lstlisting}[language=bash, caption={Menyalakan aplikasi identitas, tersedia di \berkas{sources/chapter10/aplikasi.py}}, label={lst:c10-siapkan}]
cd sources/chapter10
KUNCI_TOKEN=kunci-latihan python3 aplikasi.py &

# Keluarannya:
# Melayani di http://127.0.0.1:8070/  (Ctrl+C untuk berhenti)
\end{lstlisting}

\ifajar
  \subsection*{Latihan 1: Sesi lawan Token}
\else
  \subsection*{Prosedur 1: Sesi lawan Token}
\fi

Latihan ini membuktikan bahwa pencabutan seketika hanya dimiliki sesi, dan
bahwa isi token terbaca siapa pun yang memegangnya.

\begin{enumerate}
  \item Masuk dengan cara sesi, catat panjang penandanya dan atribut
        cookienya, seperti pada Listing \ref{lst:c10-lat1-jalan}.
  \item Pakai sesi itu untuk membuka daftar pesanan, lalu keluar, lalu
        pakai sesi yang sama sekali lagi. Catat ketiga hasilnya.
  \item Masuk dengan cara token, lalu dekode bagian tengah tokennya dengan
        \texttt{base64}. Catat apa saja yang terbaca di dalamnya.
  \item Isi Tabel \ref{tab:c10-lat1-hasil}. Tabel \ref{tab:c10-lat1-contoh}
        memperlihatkan contoh yang sudah terisi.
  \item Jawab tiga pertanyaan berikut. Pertama, mengapa token tidak dapat
        dicabut seketika? Kedua, apa yang hilang bila seluruh keterangan
        pengguna dimasukkan ke dalam token? Ketiga, kapan sesi lebih pantas
        dipakai daripada token?
\end{enumerate}

\begin{lstlisting}[language=bash, caption={Membandingkan sesi dan token, tersedia di \berkas{sources/chapter10/aplikasi.py}}, label={lst:c10-lat1-jalan}]
curl -s -D - -X POST \
  "http://127.0.0.1:8070/masuk?cara=sesi&nama=budi&sandi=budi123"
curl -s -X POST \
  "http://127.0.0.1:8070/masuk?cara=token&nama=budi&sandi=budi123"

# Keluarannya:
# Set-Cookie: sesi=4b2f8d7e...; HttpOnly; SameSite=Lax; Path=/
# {"cara": "sesi", "panjang_penanda": 32}
# {"cara": "token", "access_token": "eyJhbGciOi...", "panjang_access": 169}
\end{lstlisting}

\begin{table}[!htb]
  \centering
  \caption{Contoh hasil Latihan 1 yang sudah terisi}
  \label{tab:c10-lat1-contoh}
  \begin{tabularx}{\textwidth}{Xll}
    \hline
    \textbf{Hal} & \textbf{Sesi} & \textbf{Token} \\
    \hline
    Panjang yang dibawa klien & 32 huruf & 169 byte \\
    \hline
    Sesudah keluar & ditolak & masih berlaku \\
    \hline
  \end{tabularx}
\end{table}

\begin{table}[!htb]
  \centering
  \caption{Lembar isian hasil pengamatan Latihan 1}
  \label{tab:c10-lat1-hasil}
  \begin{tabularx}{\textwidth}{Xll}
    \hline
    \textbf{Hal} & \textbf{Sesi} & \textbf{Token} \\
    \hline
    Panjang yang dibawa klien & \dots & \dots \\
    \hline
    Sesudah keluar & \dots & \dots \\
    \hline
    Isi yang terbaca klien & \dots & \dots \\
    \hline
    Yang disimpan server & \dots & \dots \\
    \hline
  \end{tabularx}
\end{table}

\ifajar
  \subsection*{Latihan 2: Umur dan Rotasi Token}
\else
  \subsection*{Prosedur 2: Umur dan Rotasi Token}
\fi

Latihan ini membuktikan bahwa umur pendek dan rotasi menjawab dua risiko yang
berbeda.

\begin{enumerate}
  \item Masuk dengan cara token, lalu pakai \textit{access token}-nya
        segera, pada detik kedua, dan pada detik keenam. Catat ketiga
        statusnya.
  \item Ubah satu huruf terakhir token itu, lalu pakai lagi dan catat
        alasan penolakannya.
  \item Lakukan \textit{refresh} dengan \textit{refresh token}, catat pasangan baru yang
        terbit, lalu pakai ulang penanda lama seperti pada Listing
        \ref{lst:c10-lat2-jalan}.
  \item Coba penanda baru sesudah pemakaian ulang itu, lalu isi Tabel
        \ref{tab:c10-lat2-hasil}. Tabel \ref{tab:c10-lat2-contoh}
        memperlihatkan contoh yang sudah terisi.
  \item Jawab dua pertanyaan berikut. Pertama, mengapa penanda baru ikut
        tidak berlaku? Kedua, apa yang terjadi bila rotasi dipakai tanpa penangkapan
        pemakaian ulang?
\end{enumerate}

\begin{lstlisting}[language=bash, caption={Rotasi dan pemakaian ulang refresh token, tersedia di \berkas{sources/chapter10/aplikasi.py}}, label={lst:c10-lat2-jalan}]
curl -s -X POST "http://127.0.0.1:8070/refresh?refresh=PENANDA_LAMA"
curl -s -X POST "http://127.0.0.1:8070/refresh?refresh=PENANDA_LAMA"
curl -s -X POST "http://127.0.0.1:8070/refresh?refresh=PENANDA_BARU"

# Keluarannya:
# {"access_token": "eyJhbGciOi...", "refresh_token": "50dbfff1..."}
# {"pesan": "refresh token dipakai ulang", "dicabut": 2}
# {"pesan": "refresh token tidak dikenal"}
\end{lstlisting}

\begin{table}[!htb]
  \centering
  \caption{Contoh hasil Latihan 2 yang sudah terisi}
  \label{tab:c10-lat2-contoh}
  \begin{tabularx}{\textwidth}{Xll}
    \hline
    \textbf{Langkah} & \textbf{Status} & \textbf{Alasan} \\
    \hline
    Token sesudah 6 detik & 401 & token kedaluwarsa \\
    \hline
    Memakai ulang penanda lama & 401 & dipakai ulang, 2 dicabut \\
    \hline
  \end{tabularx}
\end{table}

\begin{table}[!htb]
  \centering
  \caption{Lembar isian hasil pengamatan Latihan 2}
  \label{tab:c10-lat2-hasil}
  \begin{tabularx}{\textwidth}{Xll}
    \hline
    \textbf{Langkah} & \textbf{Status} & \textbf{Alasan} \\
    \hline
    Token segera sesudah terbit & \dots & \dots \\
    \hline
    Token sesudah 2 detik & \dots & \dots \\
    \hline
    Token sesudah 6 detik & \dots & \dots \\
    \hline
    Token yang diubah satu huruf & \dots & \dots \\
    \hline
    Memakai ulang penanda lama & \dots & \dots \\
    \hline
    Penanda baru sesudah pencabutan & \dots & \dots \\
    \hline
  \end{tabularx}
\end{table}

\ifajar
  \subsection*{Latihan 3: Matriks Hak Akses}
\else
  \subsection*{Prosedur 3: Matriks Hak Akses}
\fi

Latihan ini membuktikan bahwa RBAC, ABAC, dan isolasi \textit{tenant}
terbaca sekaligus dari satu tabel.

\begin{enumerate}
  \item Jalankan alat ukurnya seperti pada Listing \ref{lst:c10-lat3-jalan},
        lalu salin hasilnya ke Tabel \ref{tab:c10-lat3-hasil}. Tabel
        \ref{tab:c10-lat3-contoh} memperlihatkan contoh yang sudah terisi.
  \item Tandai sel mana yang ditentukan peran, sel mana yang ditentukan
        kepemilikan, dan sel mana yang ditentukan \textit{tenant}.
  \item Tambahkan satu izin \texttt{baca\_laporan} pada peran staf di
        \berkas{sources/chapter10/kebijakan.py}, jalankan ulang, lalu catat
        sel mana yang berubah.
  \item Jawab tiga pertanyaan berikut. Pertama, mengapa permintaan lintas
        \textit{tenant} dijawab 404 dan bukan 403? Kedua, mengapa
        pemeriksaan \textit{tenant} ditaruh paling depan? Ketiga, apa yang
        terjadi bila penyaringan daftar dan pemeriksaan rincian memakai
        aturan yang berbeda?
\end{enumerate}

\begin{lstlisting}[language=bash, caption={Menjalankan matriks hak akses, tersedia di \berkas{sources/chapter10/ukur-akses.py}}, label={lst:c10-lat3-jalan}]
python3 ukur-akses.py

# Keluarannya:
#   Pengguna    /api/pesanan  /api/pesanan/1  /api/pesanan/3    /api/laporan
#   ani         200, 2 baris             200             404             200
#   budi        200, 2 baris             200             404             403
#   cita        200, 1 baris             200             404             403
#   dedi        200, 1 baris             404             200             403
\end{lstlisting}

\begin{table}[!htb]
  \centering
  \caption{Contoh hasil Latihan 3 yang sudah terisi}
  \label{tab:c10-lat3-contoh}
  \begin{tabularx}{\textwidth}{Xllll}
    \hline
    \textbf{Pengguna} & \textbf{Daftar} & \textbf{/1} & \textbf{/3} & \textbf{Laporan} \\
    \hline
    ani, admin kampus-a & 200, 2 baris & 200 & 404 & 200 \\
    \hline
    dedi, staf kampus-b & 200, 1 baris & 404 & 200 & 403 \\
    \hline
  \end{tabularx}
\end{table}

\begin{table}[!htb]
  \centering
  \caption{Lembar isian hasil pengukuran Latihan 3}
  \label{tab:c10-lat3-hasil}
  \begin{tabularx}{\textwidth}{Xllll}
    \hline
    \textbf{Pengguna} & \textbf{Daftar} & \textbf{/1} & \textbf{/3} & \textbf{Laporan} \\
    \hline
    ani, admin kampus-a & \dots & \dots & \dots & \dots \\
    \hline
    budi, staf kampus-a & \dots & \dots & \dots & \dots \\
    \hline
    cita, pembeli kampus-a & \dots & \dots & \dots & \dots \\
    \hline
    dedi, staf kampus-b & \dots & \dots & \dots & \dots \\
    \hline
    Sesudah staf diberi izin laporan & \dots & \dots & \dots & \dots \\
    \hline
  \end{tabularx}
\end{table}

\ifajar
  \subsection*{Latihan 4: Menambah Tenant dan Peran Baru}
\else
  \subsection*{Prosedur 4: Menambah Tenant dan Peran Baru}
\fi

Latihan ini menguji apakah aturan hak akses benar-benar terpusat, atau
sebenarnya tersebar di banyak tempat.

\begin{enumerate}
  \item Tambahkan satu \textit{tenant} ketiga beserta satu penggunanya dan
        satu pesanannya ke
        \berkas{sources/chapter10/pengguna.json}. Kata sandi pengguna baru
        disimpan sebagai \textit{salt} dan \textit{hash}, yang dibuat dengan
        perintah pada Listing \ref{lst:c10-lat4-sandi}.
  \item Tambahkan peran baru, misalnya auditor yang boleh membaca laporan
        tetapi tidak boleh membaca pesanan, ke
        \berkas{sources/chapter10/kebijakan.py}.
  \item Tambahkan kedua pengguna baru itu ke daftar pada
        \berkas{sources/chapter10/ukur-akses.py}, jalankan ulang, lalu
        catat matriksnya ke Tabel \ref{tab:c10-lat4-hasil}. Tabel
        \ref{tab:c10-lat4-contoh} memperlihatkan bentuk isian yang
        diharapkan.
  \item Catat berapa file yang harus disunting untuk kedua penambahan itu.
  \item Jawab dua pertanyaan berikut. Pertama, berapa file yang berubah, dan
        apakah jumlah itu menandakan aturannya terpusat? Kedua, apa yang
        perlu diubah bila aturannya berganti dari peran menjadi izin per
        pengguna?
\end{enumerate}

\begin{lstlisting}[language=bash, caption={Membuat \textit{salt} dan \textit{hash} kata sandi pengguna baru, tersedia di \berkas{sources/chapter10/aplikasi.py}}, label={lst:c10-lat4-sandi}]
python3 -c "import aplikasi; print(aplikasi.buat_catatan_sandi('eko123'))"

# Keluarannya, berbeda tiap kali karena salt-nya acak:
# {"salt_sandi": "aab7e943...", "hash_sandi": "97301a7f..."}
\end{lstlisting}

\begin{table}[!htb]
  \centering
  \caption{Bentuk isian Latihan 4 yang diharapkan}
  \label{tab:c10-lat4-contoh}
  \begin{tabularx}{\textwidth}{Xlll}
    \hline
    \textbf{Pengguna baru} & \textbf{Daftar} & \textbf{Rincian} & \textbf{Laporan} \\
    \hline
    auditor kampus-a & 200, 0 baris & 403 & 200 \\
    \hline
  \end{tabularx}
\end{table}

\begin{table}[!htb]
  \centering
  \caption{Lembar isian hasil pengamatan Latihan 4}
  \label{tab:c10-lat4-hasil}
  \begin{tabularx}{\textwidth}{Xlll}
    \hline
    \textbf{Pengguna baru} & \textbf{Daftar} & \textbf{Rincian} & \textbf{Laporan} \\
    \hline
    Pengguna \textit{tenant} ketiga & \dots & \dots & \dots \\
    \hline
    Peran auditor & \dots & \dots & \dots \\
    \hline
    Jumlah file yang disunting & \multicolumn{3}{l}{\dots} \\
    \hline
  \end{tabularx}
\end{table}

Setelah seluruh latihan selesai, aplikasinya dihentikan dengan Ctrl+C. Tidak
ada container yang perlu dimatikan, karena bab ini memang tidak memakainya.
